% Mesure du pouvoir calorifique de l'éthanol : une canette d'eau chauffée par une lampe à alcool (exercice 7, chapitre 10)
\begin{tikzpicture}[x=1cm,y=1cm]
\pic at (0,0) {potence={hauteur=4.6}};
% lampe à alcool
\filldraw[fill=white,draw=otInk!70,line width=0.6pt] (-0.55,0.16) rectangle (0.55,0.9);
\fill[yellow!12] (-0.5,0.2) rectangle (0.5,0.6);
\draw[otInk!70,line width=0.9pt] (0,0.9) -- (0,1.12);
\fill[otOrange] (0,1.12) .. controls (-0.18,1.25) and (-0.12,1.55) .. (0,1.75) .. controls (0.12,1.55) and (0.18,1.25) .. (0,1.12);
% canette
\filldraw[fill=otMuted!25,draw=otInk!70,line width=0.6pt,rounded corners=0.8mm] (-0.45,2.0) rectangle (0.45,3.35);
\fill[pcEau] (-0.4,2.05) rectangle (0.4,3.0);
\draw[otInk!60,line width=1pt] (-1.0,3.2) -- (-0.45,3.2);
% thermomètre
\draw[otInk!70,line width=0.8pt] (0.15,2.4) -- (0.15,4.1);
\fill[otRed!80] (0.15,2.35) circle (0.06);
\pcNote[west]{(0.45,2.6)}{(1.2,2.8)}{canette : \qty{200}{g} d'eau}
\pcNote[west]{(0.15,3.9)}{(1.2,4.0)}{thermomètre}
\pcNote[west]{(0.55,0.5)}{(1.2,0.6)}{lampe à éthanol (pesée)}
\end{tikzpicture}
\par\smallskip\legende{On mesure l'élévation de température de l'eau et la masse d'éthanol brûlé.}
